% =====================================================================
%  Thème 1 — Conversions & cercle trigonométrique — NIVEAU FACILE
%  Corrigé
% =====================================================================
\documentclass[11pt,a4paper]{article}
\usepackage[utf8]{inputenc}
\usepackage[T1]{fontenc}
\usepackage{lmodern}
\usepackage[french]{babel}
\usepackage{amsmath,amssymb,mathtools}
\usepackage{icomma}
\usepackage{xcolor}
\usepackage{colortbl}
\usepackage{tikz}
\usepackage[most]{tcolorbox}
\usepackage{enumitem}
\usepackage{array}
\usepackage[a4paper,margin=2.2cm]{geometry}
\usetikzlibrary{arrows.meta,calc}

\definecolor{primary}{RGB}{214,51,108}
\definecolor{primarydark}{RGB}{140,20,64}
\definecolor{excol}{RGB}{0,135,90}

\DeclareMathOperator{\tg}{tg}
\DeclareMathOperator{\cotg}{cotg}
\newcommand{\degr}{^{\circ}}
\newcommand{\Z}{\mathbb{Z}}

\newcounter{exo}
\newcommand{\exo}[1]{\medskip\refstepcounter{exo}%
  \noindent{\large\bfseries\color{primarydark}Exercice \theexo}%
  \ \textcolor{primary}{\textbullet}\ \textbf{#1}\par\nopagebreak\smallskip}
\newcommand{\rep}{\textcolor{excol}{$\blacktriangleright$}\ }

\setlength{\parindent}{0pt}
\pagestyle{empty}

\begin{document}

\begin{center}
{\LARGE\bfseries\color{primary}Thème 1 — Conversions \& cercle trigonométrique}\\[2pt]
{\color{primarydark}\textsc{Niveau Facile — Corrigé détaillé}}
\end{center}
\hrule height 1pt
\medskip

\exo{Des degrés vers les radians}
On multiplie par $\dfrac{\pi}{180}$ puis on simplifie.
\begin{itemize}[leftmargin=1.4em,itemsep=1pt]
\item[\textbf{(a)}] \rep $45\times\dfrac{\pi}{180}=\dfrac{45\pi}{180}=\dfrac{\pi}{4}$.
\item[\textbf{(b)}] \rep $90\times\dfrac{\pi}{180}=\dfrac{\pi}{2}$.
\item[\textbf{(c)}] \rep $150\times\dfrac{\pi}{180}=\dfrac{150\pi}{180}=\dfrac{5\pi}{6}$.
\item[\textbf{(d)}] \rep $210\times\dfrac{\pi}{180}=\dfrac{210\pi}{180}=\dfrac{7\pi}{6}$.
\item[\textbf{(e)}] \rep $300\times\dfrac{\pi}{180}=\dfrac{300\pi}{180}=\dfrac{5\pi}{3}$.
\end{itemize}

\exo{Des radians vers les degrés}
On multiplie par $\dfrac{180}{\pi}$ ; le $\pi$ se simplifie.
\begin{itemize}[leftmargin=1.4em,itemsep=1pt]
\item[\textbf{(a)}] \rep $\dfrac{\pi}{6}\times\dfrac{180}{\pi}=\dfrac{180}{6}=30\degr$.
\item[\textbf{(b)}] \rep $\dfrac{3\pi}{4}\times\dfrac{180}{\pi}=\dfrac{3\times180}{4}=135\degr$.
\item[\textbf{(c)}] \rep $\dfrac{5\pi}{3}\times\dfrac{180}{\pi}=\dfrac{5\times180}{3}=300\degr$.
\item[\textbf{(d)}] \rep $\dfrac{7\pi}{6}\times\dfrac{180}{\pi}=\dfrac{7\times180}{6}=210\degr$.
\item[\textbf{(e)}] \rep $2\pi\times\dfrac{180}{\pi}=360\degr$ (un tour complet).
\end{itemize}

\exo{Dans quel quadrant ?}
On situe l'angle par rapport à $90\degr,180\degr,270\degr,360\degr$ (soit
$\frac{\pi}{2},\pi,\frac{3\pi}{2},2\pi$).
\begin{itemize}[leftmargin=1.4em,itemsep=1pt]
\item[\textbf{(a)}] \rep $200\degr$ : entre $180\degr$ et $270\degr$ $\Rightarrow$ \textbf{quadrant III}.
\item[\textbf{(b)}] \rep $85\degr$ : entre $0\degr$ et $90\degr$ $\Rightarrow$ \textbf{quadrant I}.
\item[\textbf{(c)}] \rep $\dfrac{5\pi}{4}=225\degr$ : entre $180\degr$ et $270\degr$ $\Rightarrow$ \textbf{quadrant III}.
\item[\textbf{(d)}] \rep $\dfrac{7\pi}{4}=315\degr$ : entre $270\degr$ et $360\degr$ $\Rightarrow$ \textbf{quadrant IV}.
\item[\textbf{(e)}] \rep $\dfrac{2\pi}{3}=120\degr$ : entre $90\degr$ et $180\degr$ $\Rightarrow$ \textbf{quadrant II}.
\end{itemize}

\exo{Quel signe ?}
On lit le quadrant, puis le tableau des signes ($\cos>0$ à droite, $\sin>0$ en haut, $\tg>0$ si
$\sin$ et $\cos$ de même signe).
\begin{itemize}[leftmargin=1.4em,itemsep=1pt]
\item[\textbf{(a)}] \rep $160\degr$ : quadrant II, $\cos<0$ donc $\cos160\degr$ est \textbf{négatif}.
\item[\textbf{(b)}] \rep $250\degr$ : quadrant III, $\sin<0$ donc $\sin250\degr$ est \textbf{négatif}.
\item[\textbf{(c)}] \rep $100\degr$ : quadrant II, $\sin>0$ et $\cos<0$ donc $\tg100\degr$ est \textbf{négatif}.
\item[\textbf{(d)}] \rep $320\degr$ : quadrant IV, $\cos>0$ donc $\cos320\degr$ est \textbf{positif}.
\item[\textbf{(e)}] \rep $40\degr$ : quadrant I, $\sin>0$ donc $\sin40\degr$ est \textbf{positif}.
\end{itemize}

\exo{Valeurs remarquables du premier quadrant}
Lecture directe dans le tableau.
\begin{itemize}[leftmargin=1.4em,itemsep=1pt]
\item[\textbf{(a)}] \rep $\cos\dfrac{\pi}{3}=\dfrac12$.
\item[\textbf{(b)}] \rep $\sin\dfrac{\pi}{4}=\dfrac{\sqrt2}{2}$.
\item[\textbf{(c)}] \rep $\tg\dfrac{\pi}{3}=\sqrt3$.
\item[\textbf{(d)}] \rep $\sin\dfrac{\pi}{6}=\dfrac12$.
\item[\textbf{(e)}] \rep $\cos\dfrac{\pi}{2}=0$.
\end{itemize}

\end{document}
